\section{Direct transfer to the five-rule clause checker}
\label{sec:clausal-transfer}

The delivered checker uses the raw Eq/$D$ clause presentation, rather
than formulas with a primitive $Q$.  The newest-witness restriction must
therefore be proved directly for that presentation.  An unrestricted
native-to-clause translation can insert numerical parameters, so it
cannot simply be assumed to preserve the newest-witness condition.

\subsection{The clause interface}

Formulas are equations and $D_{x,y}B(x,y)$.  A retained clause
\[
 [A;y.B(y)]
\]
has ordinary reading $A\to\forall y B(y)$.  An ordinary assumption
$F$ is represented by $\mathsf a(F):=[0=0;y.F]$, with $y$ unused.
The five primitive families are Refl, the two-case Eval, Use, Elim,
and Bit-Intro.  Use takes a proof of the antecedent of a retained
clause and returns its body at a supplied term.  Elim opens a proved
$D_{x,y}B$ at a fresh parameter $b$ by retaining
\[
 [B(b,0);y.B(b,y)]
\]
in its branch.  Bit-Intro adds the ordinary-assumption clause
$\mathsf a(B(t_{n,\beta},\eta_{B,\beta}))$ in its step branch and
checks the conclusion $B(t_{n,\beta},Sz)$, with the selectors from
Bit-DI.  All terms are raw; no normalization is implicit.

Refl and Eval alone derive both iterator computations and the whole
equational interface.  For example, the zero case of Eval is
\[
 s=t,\quad I(0,s;u.h)=r\quad\Longrightarrow\quad r=t.
\]
Two uses, through the reflexive frame $I(0,s;u.h)$, derive the root
joining rule $s=t,\ s=r\Longrightarrow r=t$.  The successor case
applied to a one-step iterator then gives congruence in an arbitrary
term template.  All these derivations are uniform in the numerical
scope and use no introduction or elimination of $D$.  They therefore
remain valid when regenerated under any numerical renaming.

Define the closure abbreviation
\[
 q_yA(y):=D_{x,y}A(y),\qquad x\text{ unused}.
\]
For the total bit interface, case introduction of $qA$ is a direct
Bit-Intro with $\beta=1$: the witness is unused, and the empty-scope
fallback still selects cases.  Equational progress introduction of
$qE$ is direct Bit-Intro with $\beta=0$.  Closure elimination is an
Elim of $qA$, followed by Use of the resulting clause
$[A(0);y.A(y)]$.  These observations implement the $Q$ interface
without a primitive $Q$ constructor.  Hereditary premises, defined
next, supply any proofs needed below the fresh dummy witness.

\subsection{Hereditary derivations with clause contexts}

Write $K\Vdash_c F$ if every injective numerical renaming of $K,F$
into any larger finite numerical scope is derivable in the bit clause
calculus from every clause-context extension of the renamed $K$.
Renaming acts capture-avoidably on
both the antecedent and the bound body of each clause.  The extension
retains clauses, including their entire antecedent/body data; a
retained clause is not replaced by the assertion that its antecedent
or one instance of its body is true.

The cut principle needed here concerns \emph{ordinary-assumption}
clauses only:
\[
 K,\mathsf a(A)\vdash_c F,\qquad K\Vdash_c A
 \quad\Longrightarrow\quad K\vdash_c F.
 \tag{UC$_c$}
\]
It also holds when the other retained clauses are embedded in a
larger target context.  Prove it by induction on the actual first
derivation.  When Use selects the distinguished clause
$\mathsf a(A)$, replace that use by the available member of the
hereditary family for $A$; its antecedent proof of $0=0$ is no longer
needed.  Every other clause use keeps the same clause and applies
the induction hypothesis to its antecedent proof.  Elim and Intro
add their own clauses in both contexts; the replacement family is
regenerated in the enlarged numerical scope.  Thus no arbitrary
clause is cut or semantically reconstructed, and no newest-sensitive
proof is renamed.  The same argument gives cut for
$\Vdash_c$ at an ordinary-assumption clause.

Refl, Eval, and Use preserve $\Vdash_c$.  For Use the selected clause
is a member of the renamed context, and the hereditary premise
supplies its antecedent.  Elim preserves $\Vdash_c$ by extending
the source injection to send its eigenparameter to the newest
fresh target parameter; its new clause is retained literally in
the target branch.

Full formula transport is proved by induction on Eq/$D$ formulas
at the hereditary level.  The equation case is the equational
interface above.  For $D_{x,y}B(s,x,y)$, open its witness at the
newest $b$ and retain the clause
$[B(s,b,0);y.B(s,b,y)]$.  Introduce the target $D$ at $b$ with bit 1.
From its temporary ordinary assumption $B(t,b,0)$, the induction
hypothesis transports backwards to $B(s,b,0)$.  Use of the retained
clause at $Sz$ gives $B(s,b,Sz)$; forward transport gives
$B(t,b,Sz)$.  Hereditary ordinary cut composes these operations.
The result is independent of $b$, so Elim discharges it.  This proves
\[
 K\Vdash_c s=t,\quad K\Vdash_c F(s)
 \quad\Longrightarrow\quad K\Vdash_c F(t).
 \tag{HL$_c$}
\]
The induction is uniform in actual term assignments to the other
free slots.  No arbitrary term is substituted into a target proof.

\subsection{The packet and the completeness induction}

The same raw equation packet works:
\[
 T(v):=D_{x,y}\bigl(C(y,v,C(v,0,Sx))=v\bigr).
\]
At any ambient scope, $T(0)$ is obtained with bit 0 by a
progressive introduction whose successor row is a computation
theorem.  To construct $q_vT(v)$, use bit 1 on the outer dummy
witness.  In its step branch the parameter $z$ is newest; construct
$T(Sz)$ with bit 1, choosing exactly that $z$.  Its successor row
again computes to reflexivity.  This construction is repeated in
each ambient scope.  It does not assume that a direct proof of
$T(Sz)$ remains valid after $z$ ceases to be newest.

Elim followed by Use of the resulting closure clause gives
$K\Vdash_c T(t)$ for every term $t$.  Opening $T(St)$ at the
newest $b$ and using its clause at $1$ gives $Sb=St$.  The
independently equational predecessor identity gives $b=t$.
That equality has a hereditary proof in the opened clause context,
because the calculation uses only retained-clause Use and the
renaming-stable equational fragment.

Now induct on an unrestricted clause-calculus derivation, proving
hereditary derivability of its conclusion.  Refl, Eval, Use and
Elim were handled above.  A progressive introduction at witness
$0$ is directly bit 0.  For a case introduction at any term $t$,
open $T(St)$ at a newest fresh $b$.  Temporarily retain
$\mathsf a(b=t)$ and introduce the desired $D$ with bit 1.
In its step branch, hereditary transport changes the temporary
assumption $B(b,0)$ into $B(t,0)$.  Instantiate the induction
hypothesis for the original source premise in the enlarged scope:
the original source variables keep their chosen injection, the
source step parameter maps to the newest target step parameter,
and the alias parameter $b$ is skipped by this injection.
Ordinary hereditary cut supplies the zero-row assumption, and
transport backwards gives $B(b,Sz)$.

After this introduction, UC$_c$ replaces the retained ordinary
assumption $b=t$ by its hereditary proof from the opened packet.
Elim discharges $b$.  The argument applies at every numerical
renaming and clause-context extension, so it proves the hereditary
induction claim.

Conversely, every bit introduction is either ordinary cases or
equational progress.  The only extra case relative to the
baseline's witness-tested combined rule is bit 1 at an empty
scope with an equation matrix: this is case introduction at $0$.
It is derived in the unrestricted raw clause calculus by using the
nonliteral alias $w=I(0,0;u.0)$ as its witness and transporting
along $w=0$.  Since $w$ is not the literal $0$, the baseline rule
selects cases.  Thus no semantic strengthening occurs at the
empty-scope fallback.

\begin{theorem}[Five-rule bit clause calculus]
The five-rule raw Eq/$D$ clause calculus with Bit-Intro derives
exactly the same sequents from the same clause contexts as the
unrestricted clause calculus.  In particular it has the same
ordinary realization, and hence the same full first-order
bi-interpretation with iPRA.
\end{theorem}

The proof is a direct transformation into the restricted clause
calculus.  It does not rely on an intermediate proof translation
preserving which numerical parameter happens to be newest.
